% --- 1. INFORMAZIONI GENERALI ---
\section{Informazioni Generali}
\begin{itemize}
	\item \textbf{Azienda Coinvolta:} {VarGroup}
	\item \textbf{Mezzo di Comunicazione:} {Email}
	\item \textbf{Data:} 2026-03-17
\end{itemize}
Documenta una prima comunicazione conoscitiva via email tra il team SynergoUnit e i referenti dell’azienda VarGroup.
L’obiettivo della comunicazione è stato quello di manifestare interesse per il capitolato “VarGroup” e di richiedere chiarimenti, sia tecnici che organizzativi, utili allo svolgimento dello studio di fattibilità e allo sviluppo dell’MVP\textsubscript{G}.

\vspace{0.5cm}
\textbf{Referenti Azienda:}
\begin{table}[ht]
	\centering
	\renewcommand{\arraystretch}{1.2}
	\begin{tabularx}{\textwidth}{X l}
		\toprule
		\textbf{Nome e Cognome} & \textbf{Ruolo} \\
		\midrule
		{Stefano Dindo} & {Referente Principale} \\
		\bottomrule
	\end{tabularx}
\end{table}

% --- 2. CHIARIMENTI TECNICI RICHIESTI ---
\section{Chiarimenti Tecnici Richiesti}
Il gruppo ha richiesto chiarimenti sui seguenti aspetti:
\begin{enumerate}
	\item Modelli IA e Costi
	\begin{itemize}
		\item Preferenze sui modelli LLM da utilizzare (tra quelli citati nel capitolato).
		\item Modalità di accesso alle API per modelli non open source.
	\end{itemize}
	\item Testing e Coverage
	\begin{itemize}
		\item Come approcciare i test di unità per gli agenti IA
	\end{itemize}
	\item Ambito della "Remediation"
	\begin{itemize}
		\item Livello di automazione richiesto per l'MVP
	\end{itemize}
	\item Mentoring e Supporto
	\begin{itemize}
		\item Come sarà strutturato il supporto per l'acquisizione di competenze sullo stack tecnologico dell'azienda
	\end{itemize}
\end{enumerate}


% --- 3. RESOCONTO ---
\section{Resoconto}
Dalla risposta dell'azienda sono emersi i seguenti punti.

\subsection{Perimetro del progetto}
\begin{itemize}
	\item L'azienda ha precisato che la piattaforma e l'orchestratore sono già in fase di sviluppo interno. Pertanto, non sarà richiesto al gruppo di sviluppare l'intera piattaforma, bensì esclusivamente di progettare e sviluppare alcuni agenti specifici da integrare nell'ecosistema esistente.
\end{itemize}

\subsection{Modelli IA e Costi}
\begin{itemize}
	\item Gli agenti useranno LLM o SLM a seconda di cosa sarà deciso di sviluppare
\end{itemize}

\subsection{Testing e Coverage}
\begin{itemize}
	\item Il testing può prendere diverse forme:
		\begin{itemize}
			\item test per verificare che la logica delle risposte sia deterministica:
				\begin{itemize}
					\item Il routing tra agenti sia corretto: dato input X che venga invocato l'agente giusto
					\item Il parsing e la validazione delle risposte LLM: se l'LLM restituisce un JSON malformato, il sistema gestisce l'errore?
					\item Le trasformazioni di dati prima e dopo le chiamate
					\item Le policy di retry e fallback
					\item L'integrazione con strumenti esterni (es. AST parser, linter)
				\end{itemize}
			\item testare la qualità dell'output:
				\begin{itemize}
					\item Evaluation set: un campione di codice con output atteso, valutato manualmente o con un secondo LLM come "judge" (ovviamente usando SLM o LLM diversi)
					\item Assertion su struttura, non su contenuto, andremo a valutare non solo se il commento è corretto ma anche che la risposta contenga almeno un'issue con severity e line number
					\item Snapshot testing: per rilevare regressioni quando si cambia prompt o modello
				\end{itemize}
		\end{itemize}
\end{itemize}

\subsection{Ambito della "Remediation"}
\begin{itemize}
	\item È preferibile che l'agente apra direttamente la pull request sul repository di Github, ma dipende da cosa verrà richiesto
\end{itemize}

\subsection{Mentoring e Supporto}
\begin{itemize}
	\item Con gli altri gruppi è stata prevista una serie di sessioni di formazione sui temi necessari ad avere le basi per poi progettare, ricercare, studiare e fare delle proposte di soluzioni in autonomia. L'azienda rimane sempre disponibile per confronti se necessario.
\end{itemize}

% --- 4. DECISIONI FORMALI ---
\section{Decisioni Formali}
Nel caso in cui il gruppo selezionasse il presente capitolato, si intendono adottare le seguenti linee guida progettuali.

\renewcommand{\arraystretch}{1.3}
\vspace{-0.3cm}
\noindent
\begin{longtable}{c p{12cm}}
    \toprule
    \textbf{Codice} & \textbf{Decisione} \\
    \midrule
    \endfirsthead

    \toprule
    \textbf{Codice} & \textbf{Decisione} \\
    \midrule
    \endhead

    \midrule
    \multicolumn{2}{r}{\textit{Continua nella pagina successiva}} \\
    \midrule
    \endfoot

    \bottomrule
    \endlastfoot

	\textbf{VE-2.0} & Il gruppo si limiterà allo sviluppo e all'integrazione di agenti IA specifici, escludendo lo sviluppo dell'orchestratore e dell'infrastruttura di base della piattaforma. \\
    \textbf{VE-2.1} & Gli agenti useranno LLM o SLM a seconda di cosa sarà deciso di sviluppare \\
    \textbf{VE-2.2} & Testing per verificare che la logica delle risposte sia deterministica e della qualità dell'output \\
    \textbf{VE-2.3} & Raggiungere il livello di automazione richiesto dall'azienda \\
    \textbf{VE-2.4} & Seguire una serie di sessioni di formazione sui temi necessari per progettazione, ricerca, studio e proposte di soluzioni in autonomia \\
\end{longtable}


\newpage
% --- 6. FIRME ---
\section*{Approvazione e Firme}
Con le seguenti firme, i rappresentanti approvano formalmente il contenuto del presente verbale, i requisiti concordati e le decisioni prese.

\vspace{2.5cm}

\noindent
\begin{minipage}[c]{0.45\textwidth}
	\begin{center}
		\rule{6cm}{0.4pt} \\ 
		\vspace{0.2cm}
		\textbf{Per il Gruppo Synergo Unit} \\
		{Sofia De Blasi} \\
		\textit{Responsabile di Progetto}
	\end{center}
\end{minipage}
\hfill
\begin{minipage}[c]{0.45\textwidth}
	\begin{center}
		\rule{6cm}{0.4pt} \\ 
		\vspace{0.2cm}
		\textbf{Per l'Azienda {VarGroup}} \\
		{Stefano Dindo} \\
		\textit{Referente Aziendale}
	\end{center}
\end{minipage}